\documentclass[letterpaper]{article}
\usepackage{../../common/fspdoc}
\setcounter{secnumdepth}{5}
\setcounter{tocdepth}{5}
\lhead{{\em 2015 FSP Install Guide}}

% Recipe commands to run initially only in CI

% begin_fsp_run 
% fsp_validation_comment - Remove WW file update crontab entry (CI only command)
% [master]$ rm -f /etc/cron.d/wwupdatefiles
% end_fsp_run

% Include git variables
%%% This file has been generated by the vc bundle for TeX.
%%% Do not edit this file!
%%%
%%% Define Git specific macros.
\gdef\GITHash{b1e5ea39b0eb099e892994a2f56798d77c1b41ad}%
\gdef\GITAbrHash{b1e5ea3}%
\gdef\GITParentHashes{ab0141b6f1cd0193472944a1642275fb792a8f77}%
\gdef\GITAbrParentHashes{ab0141b}%
\gdef\GITAuthorName{Karl W. Schulz}%
\gdef\GITAuthorEmail{karl.w.schulz@intel.com}%
\gdef\GITAuthorDate{2015-04-17 16:52:14 -0500}%
\gdef\GITCommitterName{Karl W. Schulz}%
\gdef\GITCommitterEmail{karl.w.schulz@intel.com}%
\gdef\GITCommitterDate{2015-04-17 16:52:14 -0500}%
%%% Define generic version control macros.
\gdef\VCRevision{\GITAbrHash}%
\gdef\VCAuthor{\GITAuthorName}%
\gdef\VCDateRAW{2015-04-17}%
\gdef\VCDateISO{2015-04-17}%
\gdef\VCDateTEX{2015/04/17}%
\gdef\VCTime{16:52:14 -0500}%
\gdef\VCModifiedText{\textcolor{red}{with local modifications!}}%
%%% Assume clean working copy.
\gdef\VCModified{0}%
\gdef\VCRevisionMod{\VCRevision}%


% Define Base OS
\newcommand{\baseOS}{SLES12}
\newcommand{\LosF}{\emph{LosF}}

\begin{document}
\graphicspath{{../../common/figures/}}
\thispagestyle{empty}

%%%{\hfill\includegraphics[scale=0.14]{../figures/tron-approved}}

% Title Page

{\hspace*{4in} \includegraphics[width=1.8in]{intel_1spot_100.pdf}}

\vspace*{2cm}
\noindent {\LARGE \color{RoyalBlue} \fontfamily{phv}\selectfont 2015 Forest Peak (v15.05)} \vspace*{0.1cm} \\
\noindent {\LARGE \color{RoyalBlue} \fontfamily{phv}\selectfont Cluster Building Recipes} \\ 
\noindent\rule{6in}{0.4pt} \\

\noindent {\Large \baseOS{} Base OS} \\ \vspace{0.2cm}

\noindent {\large {\em Base Linux Edition }}

\vspace*{3in}

\noindent{\normalsize Intel Cluster Maker Team} \vspace*{0.1cm} \\
{\normalsize Copyright~{\small\copyright}~2014-2015 Intel Corporation} \vspace*{0.1cm} \\ 
{\normalsize Document Last Update: \VCDateISO} \vspace*{0.1cm} \\ 
{\normalsize Document Revision: \VCRevision} \\ \vspace*{0.1cm}

\newpage
\tableofcontents
\newpage

% Introduction  --------------------------------------------------

\section{Introduction} \label{sec:introduction}
This guide presents a simple cluster installation procedure using components
from the Forest Peak (\FSP{}) software stack. \FSP{} represents an aggregation of a number of
common ingredients required to deploy and manage an HPC Linux* cluster including
provisioning tools, resource management, I/O clients, development tools, and a variety of
scientific libraries. These packages have been pre-built with HPC integration
in mind and represent a mix of open-source components combined with \Intel{}
development and analysis tools (e.g. \Intel{} Parallel Studio XE Cluster Edition).
%This install guide assumes availablility of a {\em master} host
%that will have access to common OS repositories and the \FSP{} repository in order
%to facilitate software installs and dependency resolution.  
The documentation herein is intended to be reasonably generic, but uses the
underlying motivation of a small, 4-node cluster install to define a step-by-step
process. Several optional customizations are included and the intent is that
these collective instructions can be modified as needed for local site
customizations. 
 \\

\noindent {\bf Base Linux Edition}: this edition of the guide highlights
installation without the use of a companion configuration management system and
directly uses distro-provided package management tools for component
selection. The steps that follow also highlight specific changes to system
configuration files that are required as part of the cluster install
process. 
%Other editions of this guide provide similar install steps when using
%specific configuration management systems that can simplify the installation
%and configuration process. 
 \\

\noindent {\bf Requirements}: this recipe assumes the availability of a single head node
{\em master}, and four {\em compute} nodes. The {\em master} node is
provisioned with \baseOS{} and is subsequently configured to
provision the remaining {\em compute} nodes with Warewulf in a
stateless configuration. We assume an external \$HOME file system is
available for use on both the master and compute nodes via NFS, and that compute node BMCs are
available via IPMI from the chosen master host. \\


 \\

% -*- mode: latex; fill-column: 120; -*- 

\subsection{Inputs} \label{sec:inputs}
As this recipe details installing a cluster starting from bare-metal, there is a requirement to define IP addresses and
gather hardware MAC addresses in order to support a controlled provisioning process. These values are necessarily unique
to the hardware being used, and this document uses variable substitution (\texttt{\$\{variable\}}) in the command-line
examples that follow to highlight where local site inputs are required. A summary of the required and optional variables
used throughout this recipe are presented below. Note that while the example definitions above correspond to a small
4-node compute subsystem, the compute parameters are defined in array format to accommodate logical extension to larger
node counts. \\

\vspace*{0.2cm}
\begin{tabular}{@{}>{\textbullet}l p{7cm} l}
& \texttt{\$\{ohpc\_repo\}} & {\small \# OpenHPC repo location} \\
& \texttt{\$\{master\_name\}} & {\small \# Hostname for master server} \\
& \texttt{\$\{master\_ip\}} & {\small \# Internal IP address on master server}  \\
& \texttt{\$\{sms\_eth\_internal\}} & {\small \# Internal Ethernet interface on SMS} \\
& \texttt{\$\{eth\_provision\}} & {\small \# Provisioning interface for computes} \\
& \texttt{\$\{internal\_netmask\}} & {\small \# Subnet netmask for internal network} \\
& \texttt{\$\{ntp\_server\}} & {\small \# Local ntp server for time synchronization} \\
& \texttt{\$\{bmc\_username\}} & {\small \# BMC username for use by IPMI} \\
& \texttt{\$\{bmc\_password\}} & {\small \# BMC password for use by IPMI} \\
& \texttt{\$\{c\_ip[0]\}}, \texttt{\$\{c\_ip[1]\}}, ... & {\small \# Desired compute node addresses} \\
& \texttt{\$\{c\_bmc[0]\}}, \texttt{\$\{c\_bmc[1]\}}, ... & {\small \# BMC addresses for computes} \\
& \texttt{\$\{c\_mac[0]\}}, \texttt{\$\{c\_mac[1]\}}, ... & {\small \# MAC addresses for computes} \\
& \texttt{\$\{compute\_regex\}} & {\small \# Regex for matching compute node names (e.g. c*)} \\
\end{tabular}

\vspace*{0.2cm}
\noindent {Optional:} 
\vspace*{0.1cm}

\begin{tabular}{@{}>{\textbullet}l p{7cm} l}
& \texttt{\$\{mgs\_fs\_name\}} & {\small \# Lustre MGS mount name} \\
& \texttt{\$\{sms\_ipoib\}} & {\small \# IPoIB address for SMS server} \\
& \texttt{\$\{ipoib\_netmask\}} & {\small \# Subnet netmask for internal IPoIB} \\
& \texttt{\$\{c\_ipoib[0]\}}, \texttt{\$\{c\_ipoib[1]\}}, ... & {\small \# IPoIB addresses for computes} \\
\end{tabular}




% Base Operating System --------------------------------------------

\section{Install Base Operating System(BOS)}
In an external setting, installing the desired BOS on a {\em master} host would
typically involve booting from a DVD iso image on a new server.  However, on
the Zeus cluster, master nodes can be network booted with a minimal OS install
using an external {\em Warewulf} instance. The Warewulf configuration
provisions the OS into a ramdisk, configures local package repositories to
point to the FSP OBS server and underlying distro repos, and mounts an
NFS \$HOME file system. With this setup, all that is required to achieve a
clean BOS intall on a master host is to perform a host reboot. If, however, you
are installing a new system, insert the \baseOS{} DVD, power cycle the host,
and follow the distro provided directions to install the BOS on your chosen
{\em master} host.  


% ------------------------------------------------------------------

\section{Install FSP Components} \label{sec:basic_install}
With the BOS installed and booted, the next step is to add desired FSP packages
onto the {\em master} server in order to provide provisioning and resource
management services for the rest of the cluster. The following subsections
highlight this process.


\subsection{Enable FSP repository for local use}
To begin, enable use of the FSP repository by adding it to the local list of
available package repositories.  Note that this requires network access from
your {\em master} server to the FSP repository, or alternatively, that the FSP
repository be mirrored locally. The example below illustrates the addition
using the internal FSP checkpoint repository for the \baseOS{} distribution.


% begin_fsp_run
% fsp_validation_comment Enable FSP Repo(s)
\begin{lstlisting}[language=bash,keywords={}]
[master]$ export FSP_mirror=http://fsp-obs.pdx.intel.com:82
[master]$ zypper addrepo \
    $FSP_mirror/ForestPeak:/15.05:/Factory/SLE_12_Intel/ForestPeak:15.05:Factory.repo
[master]$ zypper -n --gpg-auto-import-keys refresh
\end{lstlisting}
% end_fsp_run

% begin_fsp_run
% % fsp_validation_comment - format space for /opt/fsp (CI only command)
% [master]$ mkfs -F -t ext4 /dev/sda1
% [master]$ mkdir -p /opt/fsp
% [master]$ mount /dev/sda1 /opt/fsp
% end_fsp_run  

\subsection{Add provisioning services on {\em master} node}
With the FSP repo enabled, we can now begin adding desired components onto the
{\em master} server. The FSP repo provides a number of aliases that group
logical components togetther in order to help aid in this process. To add
support for provisioning services, the following commands illustrate addition
of a common base package followed by the Warewulf provisioning system.


% begin_fsp_run
% fsp_validation_comment Add baseline FSP and warewulf on master node
\begin{lstlisting}[language=bash,keywords={}]
[master]$ zypper -n install -t pattern fsp-base          # adding base FSP packages
[master]$ zypper -n install -t pattern fsp-warewulf      # adding Warewulf support
\end{lstlisting}
% end_fsp_run

\subsection{Add resource management services on {\em master} node} \label{sec:add_rm}

The following command adds the SLURM workload manager server components to the
chosen {\em master} host. Note that client-side components will be added to
the corresponding compute image in a subsequent step.

% begin_fsp_run
% fsp_validation_comment Add resource manager
\begin{lstlisting}[language=bash,keywords={}]
[master]$ zypper -n install -t pattern fsp-slurm-server  # adding SLURM server support
\end{lstlisting}
% end_fsp_run

SLURM requires the delineation of the system user that runs the underlying
resource management daemons. The default configuration file that is supplied
with the FSP build of SLURM identifies this \texttt{SlurmUser} to be a
dedicated user named \texttt{slurm} and this user must exist. 
%The default location identified to save state information is
%\texttt{/var/slurm} and this directory
The following command can be used to add this user to the {\em
  master} server:

% begin_fsp_run
% fsp_validation_comment Add slurm user
\begin{lstlisting}[language=bash,keywords={}]
[master]$ useradd slurm                                  # add "slurm" user
\end{lstlisting}
% end_fsp_run

To facilitate running test jobs later on under the auspices of a resource
manager, we also add a dedicated ``test'' user as follows:

% begin_fsp_run
% fsp_validation_comment Add test user
\begin{lstlisting}[language=bash,keywords={}]
[master]$ useradd -m test                                # add "test" user
\end{lstlisting}
% end_fsp_run

\subsection{Add InfiniBand support services on {\em master} node} \label{sec:add_ofed}

The following command adds OFED support using distro-provided drivers to the
chosen {\em master} host. 

% begin_fsp_run
% fsp_validation_comment Add ofed
\begin{lstlisting}[language=bash,keywords={}]
[master]$ zypper -n install libibverbs-runtime libmlx4-rdmav2
[master]$ zypper -n install libibmad5 librdmacm1 rdma

# Load IB drivers
[master]$ service rdma start
\end{lstlisting}
% end_fsp_run

\subsection{Add Intel Cluster Checker ({\em optional})} \label{sec:add_clck}

The Intel Cluster Checker provides a suite of diagnostics that can be used to
aid in isolating hardware and software problems on an installed
cluster. Although not required, this package can be optionally added using the
command below. Note that a valid license will also need to be installed
separately in order to use this software.

% begin_fsp_run
% fsp_validation_comment Add CLCK
\begin{lstlisting}[language=bash,keywords={}]
[master]$ zypper -n install intel-clck-fsp
\end{lstlisting}
% end_fsp_run


\subsection{Complete basic Warewulf setup for {\em master} node}

At this point, all of the packages necessary to use Warewulf on the {\em
  master} host should be installed.  Next, we need to update several
configuration files in order to allow Warewulf to work with SLES and to support
local provisioning using the {\em eth1} interface.  Specific steps are as
follows:

% begin_fsp_run
% fsp_validation_comment Complete basic Warewulf setup for master node
%\begin{verbatim}

\begin{lstlisting}[language=bash,literate={-}{-}1,keywords={},upquote=true]

# Configure DHCP server to use eth1
[master]$ perl -pi -e 's/^DHCPD_INTERFACE=""/DHCPD_INTERFACE="eth1"/' /etc/sysconfig/dhcpd

# Configure Warewulf to use the default SLES tftp location
[master]$ perl -pi -e "s#\#tftpdir = /var/lib/#tftpdir = /srv/#" /etc/warewulf/provision.conf

# Enable tftp service
[master]$ perl -pi -e "s/^\s+disable\s+= yes/ disable = no/" /etc/xinetd.d/tftp

# Update Warewulf http config to use the SLES version of mod_perl
[master]$ export MODFILE=/etc/apache2/conf.d/warewulf-httpd.conf
[master]$ perl -pi -e "s#modules/mod_perl.so\$#/usr/lib64/apache2/mod_perl.so#" $MODFILE

# Enable http access for Warewulf cgi-bin directory to support newer apache syntax
[master]$ perl -pi -e "s/cgi-bin>\$/cgi-bin>\n Require all granted/" $MODFILE
[master]$ perl -pi -e "s/Allow from all/Require all granted/" $MODFILE
[master]$ perl -ni -e "print unless /^\s+Order allow,deny/" $MODFILE

# Define eth1 interface for provisioning
[master]$ ifconfig eth1 <master_ip> netmask 255.255.255.0 up

# Restart relevant services to support provisioning
[master]$ service xinetd restart                         # tftp services
[master]$ service mysql restart                          # mysql database
[master]$ service apache2 restart                        # web server
\end{lstlisting}
%\end{verbatim}
% end_fsp_run

%%% The steps above included starting the {\em eth1} interface on the {\em master}
%%% host in order to provide provisioning services locally for the cluster. In this
%%% install recipe, the {\em master} host will also provide resource manager
%%% services and requires a hostname to be associated with the {\em eth1} interface
%%% so that the compute hosts can reference it by name. Portions of the
%%% \texttt{/etc/hosts} file on the {\em master} host will be automatically
%%% propagated to the provisioned hosts and a hostname entry that maps to this {\em
%%%   eth1} interface can be added as follows:
%%% 
%%% % abegin_fsp_run
%%% \begin{lstlisting}[language=bash,literate={-}{-}1,keywords={},upquote=true]
%%% # add local hostname entry for eth1 
%%% [master]$ echo "<master_ip> <master_hostname>-eth1" >> /etc/hosts   
%%% 
%%% # import (sync) the updated hosts file with Warewulf
%%% [master]$ wwsh file sync dynamic_hosts
%%% \end{lstlisting}
%%% % aend_fsp_run

\subsection{Define {\em compute} image for provisioning}

With the provisioning services enabled, the next step is to define and
customize a system image that can subsequently be used to provision one or more
{\em compute} nodes. The following subsections highlight this process.

\subsubsection{Build initial BOS image}

The FSP build of Warewulf includes specific enhancements enabling support for
SLES~12. The following steps illustrate the process to build a minimal, default
image for use with Warewulf. In this example, the resulting image will be
stored locally on the {\em master} host in the
\texttt{/opt/fsp/admin/images/sles12} directory.

% begin_fsp_run
% fsp_validation_comment Create compute image for Warewulf
\begin{lstlisting}[language=bash,literate={-}{-}1,keywords={},upquote=true]
# Define chroot location for SLES 
[master]$ export CHROOT=/opt/fsp/admin/images/sles12

# Configure additional required drivers for bootstrap
[master]$ export WW_CONF=/etc/warewulf/bootstrap.conf
[master]$ echo "drivers += af_packet, dns_resolver, auth_rpcgss, lockd, sunrpc" >> $WW_CONF
[master]$ echo "drivers += dm_mod, dcdbas, dell_rbu, ipmi_msghandler, shpchp, ehci_pci" >> $WW_CONF
[master]$ echo "drivers += wmi, processor, button, sg, igb, i2c_algo_bit, dca, ptp" >> $WW_CONF
[master]$ echo "drivers += pps_core, ahci, libahci, libata, scsi_mod" >> $WW_CONF
[master]$ echo "drivers += nfs_acl, nfsv3, nfsv4, coretemp" >> $WW_CONF


# Build bootstrap kernel and initial base image
[master]$ wwbootstrap `uname -r`              # create warewulf bootstrap image
[master]$ mkdir -p --mode=755 $CHROOT         # create chroot housing dir
[master]$ mkdir -m 755 $CHROOT/dev            # create chroot /dev dir
[master]$ mknod -m 666 $CHROOT/dev/zero c 1 5 # create /dev/zero device
[master]$ wwmkchroot sles-12 $CHROOT          # create base image

\end{lstlisting}
% end_fsp_run

\subsubsection{Add FSP components}

The \texttt{wwmkchroot} process used in the previous step is designed to
provide a minimal SLES configuration. Next, we add additional components to
include resource management client services, InfiniBand drivers, and other
additional packages to support the default FSP environment.  This process uses
the \texttt{chroot} command to augment the base provisioning image and will
access the BOS and FSP repositories to resolve package install requests. To
access the remote repositories by hostname (and not IP addresses), the chroot
environment needs to be updated to enable DNS resolution. Assuming that
the {\em master} host has a working DNS configuration in place, the chroot environment can
be updated with a copy of the configuration as follows:

% begin_fsp_run
\begin{lstlisting}[language=bash,literate={-}{-}1,keywords={},upquote=true]
[master]$ cp -p /etc/resolv.conf $CHROOT/etc/resolv.conf
\end{lstlisting}
% end_fsp_run


% begin_fsp_run
% fsp_validation_comment Add FSP components to compute instance
\begin{lstlisting}[language=bash,literate={-}{-}1,keywords={},upquote=true]
# Add SLURM client support
[master]$ chroot $CHROOT zypper -n install -t pattern fsp-slurm-client

# Add IB support
[master]$ chroot $CHROOT zypper -n install libibverbs-runtime libmlx4-rdmav2
[master]$ chroot $CHROOT zypper -n install libibmad5 librdmacm1 rdma

# Add Network Time Protocol (NTP) support
[master]$ chroot $CHROOT zypper -n install ntp

# Include FSP default modules user environment
[master]$ chroot $CHROOT zypper -n install FSP-lmod-defaults-intel

# Enable ssh access 
[master]$ chroot $CHROOT systemctl enable sshd.service

# Remove default hostname to allow WW to provision network names
[master]$ mv $CHROOT/etc/hostname $CHROOT/etc/hostname.orig
\end{lstlisting}
% end_fsp_run

% begin_fsp_run 
% fsp_validation_comment - Restart openibd (CI only command)
% [master]$ service openibd restart
% end_fsp_run

\subsubsection{Customize system configuration} \label{sec:master_customization}

Prior to assembling the image, it is advantageous to perform any additional
customizations within the chroot environment created for the desired {\em
  compute} instance. The following steps document the process to add a local
{\em ssh} key created by Warewulf to support remote access, identify the
resource manager server, and enable NFS mounting of a \$HOME file system and
the public FSP install path (\texttt{/opt/fsp/pub}) that will be hosted by the
{\em master} host in this example configuration.  The NFS exporting options use
an address/netmask combination to limit the export scope to the defined compute nodes.

%%Note that the example shown
%%here leverages the {\em eth1} hostname identified for the {\em master} host and
%%specifically exports the desired NFS paths to each compute node
%%individually. In practice, however, a single entry would likely be used to
%%export the mount point to a private subnet suitably sized to comprise all
%%desired endpoints.

% begin_fsp_run
% fsp_validation_comment Add ssh key and enable NFS mount of /opt/fsp/pub
\begin{lstlisting}[language=bash,literate={-}{-}1,keywords={},upquote=true]
# add new cluster key to base image
[master]$ cat ~/.ssh/cluster.pub >> $CHROOT/root/.ssh/authorized_keys

# add NFS client mounts of /home and /opt/fsp/pub to base image
[master]$ echo "<master_ip>:/home /home nfs nfsvers=3,rsize=1024,wsize=1024,cto 0 0" >> \
     $CHROOT/etc/fstab
[master]$ echo "<master_ip>:/opt/fsp/pub /opt/fsp/pub nfs nfsvers=3 0 0" >> \
     $CHROOT/etc/fstab

# Identify resource manager server IP in SLURM config file on computes
[master]$ perl -pi -e "s/ControlMachine=\S+/ControlMachine=<master_ip>/" \
    $CHROOT/etc/slurm/slurm.conf

# Identify resource manager hostname on master host
[master]$ perl -pi -e "s/ControlMachine=\S+/ControlMachine=<master_hostname>/" \
    /etc/slurm/slurm.conf

# Export /home and FSP public packages from master server to cluster compute nodes
[master]$ echo "/home <internal_netmask>(rw,no_subtree_check,fsid=10,no_root_squash)" \
     >> /etc/exports
[master]$ echo "/opt/fsp/pub <internal_netmask>(ro,no_subtree_check,fsid=11)" \
     >> /etc/exports
[master]$ exportfs -a
[master]$ service nfsserver restart
\end{lstlisting}
% end_fsp_run


\subsubsection{Additional Customizations ({\em optional})}

This section highlights common additional customizations that
can {\em optionally} be applied to the
local cluster environment. These customizations include:

\begin{itemize*}
\item Increase memlimts to support InfiniBand
\item Restrict ssh access to compute resources
\item Add Lustre client
\end{itemize*}

\noindent Details on the steps required for each of these customizations is
discussed further in the following sections.

\paragraph{Increase locked memory limits}

In order to utilize \InfiniBand{} or Omni-Path as the underlying high speed interconnect, it is
generally necessary to increase the locked memory settings for system
users. This can be accomplished by updating
the \texttt{/etc/security/limits.conf} file and this should be performed within
the {{\em compute} image and on all job submission hosts. In this recipe, jobs
are submitted from the {\em master} host, and the following commands can be
used to update the maximum locked memory settings on both the master host and
the compute image:

% begin_ohpc_run
% ohpc_validation_newline 
% ohpc_validation_comment Update memlock settings
\begin{lstlisting}[language=bash,keywords={},upquote=true]
# Update memlock settings on master
[sms](*\#*) perl -pi -e 's/# End of file/\* soft memlock unlimited\n$&/s' /etc/security/limits.conf
[sms](*\#*) perl -pi -e 's/# End of file/\* hard memlock unlimited\n$&/s' /etc/security/limits.conf

# Update memlock settings within compute image
[sms](*\#*) perl -pi -e 's/# End of file/\* soft memlock unlimited\n$&/s' $CHROOT/etc/security/limits.conf
[sms](*\#*) perl -pi -e 's/# End of file/\* hard memlock unlimited\n$&/s' $CHROOT/etc/security/limits.conf
\end{lstlisting}
% end_ohpc_run


 

% begin_fsp_run
% fsp_validation_comment Update memlock on master
\begin{lstlisting}[language=bash,keywords={},upquote=true]
# Update memlock settings on master
[master]$ echo "* soft memlock unlimited" >> /etc/security/limits.conf
[master]$ echo "* hard memlock unlimited" >> /etc/security/limits.conf

# Update memlock settings within compute image
[master]$ echo "* soft memlock unlimited" >> $CHROOT/etc/security/limits.conf
[master]$ echo "* hard memlock unlimited" >> $CHROOT/etc/security/limits.conf
\end{lstlisting}
% end_fsp_run

\paragraph{Enable ssh control via resource manager} 

An additional optional customization that is recommended is to
restrict \texttt{ssh} access on compute nodes to only allow access by users who
have an active job associated with the node. This can be enabled via the use of
a pluggable authentication module (PAM) provided as part of the SLURM package
installs. To enable this feature within the {\em compute} image, issue the following:




% begin_fsp_run
% fsp_validation_comment Enbale slurm pam module
\begin{lstlisting}[language=bash,keywords={},upquote=true]
[master]$ echo "account    required     pam_slurm.so" >> $CHROOT/etc/pam.d/sshd
\end{lstlisting}
% end_fsp_run

\paragraph{Add Lustre client} 

To add \Lustre{} client support on the cluster, it necessary to install the client
and associated modules on each host needing to access a \Lustre{} file system. In
this recipe, it is assumed that the \Lustre{} file system is hosted by servers
that are pre-existing and are not part of the install process. Outlining the
variety of \Lustre{} client mounting options is beyond the scope of this document,
%(please consult \Lustre{} documentation for more details on failover configuration
%support and networking options), 
but the general requirement is to add a mount entry for the desired file system
that defines the management server (MGS) and underlying network transport
protocol. To add client mounts on both the {\em master} server and {\em
compute} image, the following commands can be used. Note that the \Lustre{} file
system to be mounted is identified by the \texttt{\$\{mgs\_fs\_name\}} variable. 
In this example, the file system is configured to be mounted locally
as \path{/mnt/lustre}.

% SLES-specific addition
 Additionally, note that your default SLES environment
  may not allow loading of the necessary Lustre kernel modules. Consequently,
  the example below includes steps which update the {\texttt
    /etc/modprobe.d/10-unsupported-modules.conf} file to allow loading of the necessary
  modules.

% begin_fsp_run
% fsp_validation_comment Enbale slurm pam module
\begin{lstlisting}[language=bash,keywords={},upquote=true]
# Add Lustre client software to master host
[master]$ zypper -n install lustre-client-fsp lustre-client-fsp-modules
# Update config to allow Lustre modules to be loaded on master host
[master]$ perl -pi -e "s/^allow_unsupported_modules 0/allow_unsupported_modules 1/" \
     /etc/modprobe.d/10-unsupported-modules.conf
# Mount the Lustre file system on master host
[master]$ mkdir /mnt/lustre
[master]$ mount -t lustre <mgs_node>:/lustrewt /mnt/lustre
\end{lstlisting}
% end_fsp_run

% begin_fsp_run
% fsp_validation_comment Enable lustre in WW compute image
\begin{lstlisting}[language=bash,keywords={},upquote=true]
# Include Lustre client software in compute image
[master]$ zypper --root=$CHROOT -n install lustre-client-fsp lustre-client-fsp-modules
# Update config to allow Lustre modules to be loaded on compute hosts
[master]$ perl -pi -e "s/^allow_unsupported_modules 0/allow_unsupported_modules 1/" \
     $CHROOT/etc/modprobe.d/10-unsupported-modules.conf
# Include mount point and file system mount in compute image
[master]$ mkdir $CHROOT/mnt/lustre
[master]$ echo "mount -t lustre <mgs_node>:/lustrewt /mnt/lustre" >> $CHROOT/etc/fstab
\end{lstlisting}
% end_fsp_run

\paragraph{Enable ORCM RAS subsystem} 

FSP includes a monitoring capability that can optionally be installed to
support unified data collection across the cluster. This monitoring system is
based on ORCM, the open resilient cluster manager project, and provides a
convenient way to aggregate key node metrics including environmental sensors,
processor power, and resource utilization into a centralized database for
historical logging and analysis.  In a typical HPC configuration, ORCM data
collection relies on having IPMI access to the baseboard management controller
(BMC) on each back-end compute resource. The data collection itself can either
be performed using {\em in-band} methods where a daemon is instantiated on each
compute resource or via {\em out-of-band} methods using external IPMI access on
a management network. To include all the components necessary to enable an ORCM
server running on the {\em master} host, issue the following:

% begin_fsp_run
% fsp_validation_comment Install ORCM server
\begin{lstlisting}[language=bash]
[master]$ zypper -n install -t pattern fsp-orcm-server
# Set master host as data aggregator
[master]$ perl -pi -e "s/aggregator_hostname/localhost/" /etc/orcm-site.xml
\end{lstlisting}
% end_fsp_run

To use the {\em in-band} method of collecting data from compute resources, the
ORCM client daemon is required to be available on each compute node. This can
be included in the compute image as follows:

% begin_fsp_run
% fsp_validation_comment Install ORCM client
\begin{lstlisting}[language=bash]
[master]$ zypper --root=$CHROOT -n install -t pattern fsp-orcm-client
# Set master host as data aggregator
[master]$ perl -pi -e "s/aggregator_hostname/<master_ip>/" /etc/orcm-site.xml
\end{lstlisting}
% end_fsp_run

\subsubsection{Import files} \label{sec:file_import}

The Warewulf system includes functionality to import arbitrary files from the
provisioning server for distribution to managed hosts. This is one way
to distribute user credentials across to {\em compute} hosts. To
import local file-based credentials, issue the following:

% begin_fsp_run
% fsp_validation_comment Import credentials
\begin{lstlisting}[language=bash,literate={-}{-}1,keywords={},upquote=true]
[master]$ wwsh file import /etc/passwd                                                                                                       
[master]$ wwsh file import /etc/group
[master]$ wwsh file import /etc/shadow 
\end{lstlisting}
% \end_fsp_run

Similarly, to import the cryptographic key that is required by the {\em munge}
authentication library to be available on every host in the resource management
pool, issue the following:

% begin_fsp_run
\begin{lstlisting}[language=bash,literate={-}{-}1,keywords={},upquote=true]
[master]$ wwsh file import /etc/munge/munge.key
\end{lstlisting}
% \end_fsp_run


\subsubsection{Assemble vNFS image}

With the local site customizations in place, the following step uses the
\texttt{wwvnfs} command to assemble a vNFS image from the chroot environment
defined for the {\em compute} instance. 

% begin_fsp_run
% fsp_validation_comment Assemble vNFS
\begin{lstlisting}[language=bash,literate={-}{-}1,keywords={},upquote=true]
[master]$ wwvnfs -y --chroot $CHROOT
\end{lstlisting}
% end_fsp_run

\subsubsection{Register nodes for provisioning}

In preparation for provisioning, we can now define the desired network settings
for four example compute nodes with the underlying provisioning system and
restart the \texttt{dhcp} service. Note the use of variable names for the
desired compute node IP and MAC addresses which should be modified to
accommodate local settings and hardware. The final step in this process
associates the vNFS image assembled in previous steps with the newly defined
compute nodes, utilizing the user credential files and munge key that were
imported in \S\ref{sec:file_import}.

% begin_fsp_run
% fsp_validation_comment Add hosts to cluster

\begin{lstlisting}[language=bash,keywords={},upquote=true,basicstyle=\footnotesize\ttfamily]
# Define four compute nodes and network settings 
[master]$ wwsh -y node new c1 --ipaddr=<c1_ip> --hwaddr=<c1_mac> 
[master]$ wwsh -y node new c2 --ipaddr=<c2_ip> --hwaddr=<c2_mac> 
[master]$ wwsh -y node new c3 --ipaddr=<c3_ip> --hwaddr=<c3_mac> 
[master]$ wwsh -y node new c4 --ipaddr=<c4_ip> --hwaddr=<c4_mac> 

# Define provisioning image for hosts
[master]$ wwsh -y provision set c[1-4] --vnfs=sles12 --bootstrap=`uname -r` \
     --files=dynamic_hosts,passwd,group,shadow,munge.key 

# Restart dhcp / update PXE
[master]$ service dhcpd restart
[master]$ wwsh pxe update
\end{lstlisting}
% end_fsp_run

\subsection{Boot compute nodes}

At this point, the {\em master} server should be able to boot the newly defined
compute nodes. Assuming that the compute node BIOS settings are configured to
boot over PXE, all that is required to initiate the provisioning process is to
power cycle each of the desired hosts using IPMI access.
The following commands use the \texttt{ipmitool} utility to initiate power
resets on each of the four compute hosts. Note that the utility requires that
the \texttt{IPMI\_PASSWORD} environment variable be set with the local BMC password in
order to work interactively.

% begin_ohpc_run
% ohpc_comment_header Boot compute nodes \ref{sec:boot_computes}
\begin{lstlisting}[language=bash,keywords={},upquote=true]
[sms](*\#*) for ((i=0; i<${num_computes}; i++)) ; do
              ipmitool -E -I lanplus -H ${c_bmc[$i]} -U ${bmc_username} chassis power reset
        done
\end{lstlisting} 
% end_ohpc_run

Once kicked off, the boot process should take less than 5 minutes (depending on
BIOS post times) and you can verify that the compute hosts are available via
ssh, or via parallel ssh tools to multiple hosts. For example, to run a command
on the newly imaged compute hosts using \texttt{pdsh}, execute the following:

%\iftoggle{isCentOS}{\clearpage}
  
\begin{lstlisting}[language=bash]
[sms](*\#*) pdsh -w c[1-4] uptime
c1  05:03am  up   0:02,  0 users,  load average: 0.20, 0.13, 0.05
c2  05:03am  up   0:02,  0 users,  load average: 0.20, 0.14, 0.06
c3  05:03am  up   0:02,  0 users,  load average: 0.19, 0.15, 0.06
c4  05:03am  up   0:02,  0 users,  load average: 0.15, 0.12, 0.05
\end{lstlisting}
 
The following commands use the \texttt{ipmitool} utility to initiate power
resets on each of the four compute hosts. Note that the utility requires that
the \texttt{IPMI\_PASSWORD} environment variable be set with the local BMC password in
order to work interactively.

% begin_fsp_run
% fsp_validation_comment Boot compute nodes

\begin{lstlisting}[language=bash,keywords={},upquote=true]
[master]$ ipmitool -E -I lanplus -H <c1_bmc> -U root chassis power reset   # power cycle c1
[master]$ ipmitool -E -I lanplus -H <c2_bmc> -U root chassis power reset   # power cycle c2
[master]$ ipmitool -E -I lanplus -H <c3_bmc> -U root chassis power reset   # power cycle c3
[master]$ ipmitool -E -I lanplus -H <c4_bmc> -U root chassis power reset   # power cycle c4
\end{lstlisting} 

% end_fsp_run

Once kicked off, the boot process should take less than 2 minutes and you can
verify that the compute hosts are available via ssh, or via parallel ssh tools to multiple
hosts. For example, to wait 120 seconds, and run a command on the newly imaged
compute hosts using \texttt{pdsh}, execute the following:

% begin_fsp_run
\begin{lstlisting}[language=bash]
[master]$ sleep 120
[master]$ /opt/fsp/admin/pdsh/bin/pdsh -w c[1-4] uptime
c1  05:03am  up   0:02,  0 users,  load average: 0.20, 0.13, 0.05
c2  05:03am  up   0:02,  0 users,  load average: 0.20, 0.14, 0.06
c3  05:03am  up   0:02,  0 users,  load average: 0.19, 0.15, 0.06
c4  05:03am  up   0:02,  0 users,  load average: 0.15, 0.12, 0.05
\end{lstlisting}
% end_fsp_run



\section{Install FSP Development Components}

The install procedure outlined in \S\ref{sec:basic_install}
highlighted the steps necessary to install a {\em master} host,
assemble and customize a {\em compute} image, and provision several
compute hosts from bare-metal.  With these steps completed, 
%a {\em compute} Once the completion of the
%basic install procedure outlined in Section~\ref{sec:basic_install} is
%complete, 
additional FSP-provided packages can now be added to support a flexible HPC
development environment including development tools, C/C++/Fortran compilers,
MPI stacks, and a variety of 3rd party libraries. The following subsections
highlight the additional software installation procedures, including the
addition of Intel licensed software (e.g. Composer compiler suite, Intel
MPI). It is assumed that the end-site administrator will procure and install
the necessary licenses in order to use the Intel proprietary software.

\subsection{Development Tools}

To aid in general development efforts, FSP provides recent versions of the GNU
autotools collection and the {\em valgrind} memory debugger. These can be installed as follows:

% begin_fsp_run
% fsp_validation_comment Install dev tools
\begin{lstlisting}[language=bash,keywords={},literate={-}{-}1]
[master]$ zypper -n install -t pattern fsp-autotools
[master]$ zypper -n install valgrind-fsp
\end{lstlisting}
% end_fsp_run

\subsection{Compilers}

FSP presently provides two compiler families ({GNU} and {Intel}) that are
integrated within the underlying modules-environment system in a hierarchical
fashion. End users of a FSP system can choose to access one compiler at a time
and will be presented with additional compiler-dependent software as a function
of which compiler toolchain is currently loaded. Each compiler toolchain can be
installed separately and the following commands illustrate the installation of
both along with any necessary dependencies:

% begin_fsp_run
% fsp_validation_comment Install compilers
\begin{lstlisting}[language=bash]
[master]$ zypper -n install gnu-compilers-fsp intel-compilers-fsp
\end{lstlisting}
% end_fsp_run

\subsection{Performance Tools}

To aid in application performance analysis, FSP provides a variety of
open-source and Intel licensed software. These can be installed as follows:

% begin_fsp_run
% fsp_validation_comment Install performance tools
\begin{lstlisting}[language=bash,keywords={},literate={-}{-}1]
[master]$ zypper -n install papi-fsp
[master]$ zypper -n install intel-itac-fsp
[master]$ zypper -n install intel-vtune-fsp
[master]$ zypper -n install intel-advisor-fsp
[master]$ zypper -n install intel-inspector-fsp
\end{lstlisting}
% end_fsp_run

\subsection{MPI Stacks} \label{sec:mpi}

For MPI development support, FSP presently provides pre-packaged builds for
three MPI families: 

\begin{itemize*}
\item Intel~MPI
\item OpenMPI
\item MVAPICH2
\end{itemize*}
 For ABI consistency, each of the open-source MPI families (OpenMPI and
 MVAPICH2) is built against each of the two supported compiler families
 resulting in total of four build combinations.  The Intel MPI stack is also
 configured to support both the GNU and Intel compiler toolchain directly, but
 is packaged as a single RPM. Installation of all of the MPI family instances,
 can be accomplished via the following command. Note the use of wildcards
 (\texttt{*}) in this example in order to install both GNU and Intel builds for
 OpenMPI and MVAPICH2.

% begin_fsp_run
% fsp_validation_comment Install MPI
\begin{lstlisting}[language=bash]
[master]$ zypper -n install  openmpi-*-fsp mvapich2-*-fsp intel-mpi-fsp
[master]$ zypper -n install -t pattern fsp-imb
\end{lstlisting}
% end_fsp_run

\subsection{Install Licenses for Intel Software}

For internal testing convenience, an RPM containing valid licenses for Intel
compilers and MPI is available. To install, issue:

% begin_fsp_run
\begin{lstlisting}[language=bash,keywords={},upquote=true]
[master]$ zypper -n install intel_licenses
\end{lstlisting}
% end_fsp_run

\subsection{3rd Party Libraries and Tools}

FSP provides pre-packaged builds for a number of popular open-source
tools and libraries used by HPC applications and developers. For
example, FSP provides builds for FFTW and HDF5 (including serial and parallel
I/O support), and the GNU Scientific Library (GSL). Again, multiple builds of
each package are available in the FSP repository to support multiple compiler
and MPI family combinations where appropriate. The general naming convention
for builds provided by FSP is to append the compiler and MPI family name that
the library was built against directly into the package name. For example,
libraries that do not require MPI as part of the build process adopt the
following RPM name: \\

\noindent
\texttt{package-<compiler\_family>-fsp-<package\_version>-<release>.rpm} \\

\noindent Packages that require MPI as part of the build expand upon this convention to
additionally include the MPI family name as follows: \\

\noindent
\texttt{package-<compiler\_family>-<mpi\_family>-fsp-<package\_version>-<release>.rpm} \\

To illustrate this further, the command below queries the locally configured
repositories to identify all of the available FFTW packages that were built
with the GNU toolchain. The resulting output that is included shows that
pre-built versions are available for each of the supported MPI families
presented in \S\ref{sec:mpi}.

\begin{lstlisting}[language=bash]
[master]$ zypper search -t package fftw-gnu-*-fsp
Loading repository data...
Reading installed packages...

S | Name                  | Summary                          | Type   
--+-----------------------+----------------------------------+--------
  | fftw-gnu-impi-fsp     | A Fast Fourier Transform library | package
  | fftw-gnu-mvapich2-fsp | A Fast Fourier Transform library | package
  | fftw-gnu-openmpi-fsp  | A Fast Fourier Transform library | package
\end{lstlisting}

Note that FSP provided 3rd party builds are configured to be installed
into a common top-level repository so that they can be easily exported to
desired hosts within the cluster. This common top-level path
(\texttt{/opt/fsp/pub}) was previously configured to be mounted on {\em
  compute} nodes in \S\ref{sec:master_customization}, so the packages will be
immediately available for use on the cluster after installation on the {\em
  master} host.  For convenience, FSP provides package aliases for these 3rd
party libraries and utilities that can be used to install all of the available
compiler/MPI family permutations. To install all of the available package
offerings within FSP, issue the following:

% begin_fsp_run
% fsp_validation_comment Install 3rd party libraries
\begin{lstlisting}[language=bash,keywords={},upquote=true]
[master]$ zypper -n install -t pattern fsp-numpy    # adds 2 numerical python packages
[master]$ zypper -n install -t pattern fsp-scipy    # adds 6 scientific python packages
[master]$ zypper -n install -t pattern fsp-hdf5     # adds 2 HDF5 packages
[master]$ zypper -n install -t pattern fsp-netcdf   # adds 6 NetCDF packages
[master]$ zypper -n install -t pattern fsp-gsl      # adds 2 GSL packages
[master]$ zypper -n install -t pattern fsp-metis    # adds 2 METIS packages
[master]$ zypper -n install -t pattern fsp-fftw     # adds 6 FFTW packages
[master]$ zypper -n install -t pattern fsp-phdf5    # adds 6 (parallel) HDF5 packages
[master]$ zypper -n install -t pattern fsp-petsc    # adds 6 PETSC packages
[master]$ zypper -n install -t pattern fsp-boost    # adds 6 Boost packages
\end{lstlisting}
% end_fsp_run

\section{Resource Manager Startup}

In section \S\ref{sec:basic_install}, the SLURM resource manager was installed
and configured for use on both the {\em master} host and {\em compute} node
instances. With the cluster nodes up and functional, we can now startup the
resource manager services in preparation for running user jobs. Generally, this
is a two-step process that requires starting up the controller daemons on the {\em
  master} host and the client daemons on each of the {\em compute} hosts.  
%Since the {\em compute} hosts were booted into an image that included the SLURM client
%components, the daemons should already be running on the {\em compute}
%hosts. 
Note that SLURM leverages the use of the {\em munge} library to provide
authentication services and this daemon also needs to be running on all hosts
requiring within the resource management pool. 
%The munge daemons should already
%be running on the {\em compute} subsystem at this point, 
The following commands can be used to startup the necessary services to support
resource management under SLURM.
%,  onso the steps required
%to startup the required services on the {\em master} host are as follows: 

% begin_fsp_run
% fsp_validation_comment Startup RM
\begin{lstlisting}[language=bash]
# start munge and slurm controller on master host
[master]$ service munge start
[master]$ service slurmctld start

#  start munge and slurm clients on compute hosts
[master]$ /opt/fsp/admin/pdsh/bin/pdsh -w c[1-4] service munge start
[master]$ /opt/fsp/admin/pdsh/bin/pdsh -w c[1-4] service slurmd start
\end{lstlisting}
% end_fsp_run

In the default configuration, the {\em compute} hosts will be initialized in an
{\em unkown} state. To place the hosts into production such that they are
eligble to schedule user jobs, issue the following:

% begin_fsp_run
% fsp_validation_comment - Sleep for 5 seconds (CI only command)
% [master]$ sleep 5
% end_fsp_run

% begin_fsp_run
\begin{lstlisting}[language=bash]
[master]$ scontrol update nodename=c[1-4] state=idle
\end{lstlisting}
% end_fsp_run


\section{Run a Test Job}

With the resource manager enabled for production usage, users should now be
able to run jobs. 
%The Zeus system includes a ``test'' user account housed on the
%NFS file system that was propagated to the {\em compute} nodes in this recipe
%via the import of user credentials in \S\ref{sec:file_import}. 
Recall that we added a ``test'' user on the {\em master} host in
\S\ref{sec:add_rm} that can now be used to run an example test job.  FSP
includes a simple ``hello-world'' MPI application in the
\texttt{/opt/fsp/pub/examples} directory that can be used for this quick
compilation and execution.  To use the test account to compile and execute the 
application interactively through the resource manager, execute the following:

\begin{lstlisting}[language=bash,keywords={}]
# switch to "test" user
[master]$ su - test

# Compile and execute with resource manager
[test@master ~]$ mpicc -o hello -O3 /opt/fsp/pub/examples/mpi/hello.c

[test@master ~]$ srun -n 8 -N 2 ./hello

Hello, world (8 procs total)
    --> Process #   4 of   8 is alive. -> c2
    --> Process #   0 of   8 is alive. -> c1
    --> Process #   5 of   8 is alive. -> c2
    --> Process #   1 of   8 is alive. -> c1
    --> Process #   6 of   8 is alive. -> c2
    --> Process #   2 of   8 is alive. -> c1
    --> Process #   7 of   8 is alive. -> c2
    --> Process #   3 of   8 is alive. -> c1
\end{lstlisting}

\end{document}

